\section{Periodic-orbit families}
\label{sec:periodic}

The centre modes identified in Section~\ref{subsec:centrefreq} provide
local seeds for Lyapunov families of periodic orbits. We compute ten such
families, chosen to sample doubly-elliptic equilibria, saddle$\times$centre
equilibria, and selected low-order near-commensurability regions. This
section describes the shooting formulation and continuation procedure,
classifies the corrected orbits by their Floquet multipliers, and examines
the approach of two families to the period-doubling threshold. The computed
families are summarized in Table~\ref{tab:periodic_families}.

% AUTO-GENERATED by paper/make_tables.py from results/*.csv. Do not edit by hand.
\begin{table}[t]
\centering
\caption{Computed periodic-orbit families. Each family is a Lyapunov family seeded from a simple centre mode of a parent equilibrium and continued in amplitude; the Lyapunov centre theorem applies at every parent (the exact non-resonance condition holds, including at the saddle$\times$centre parents of families 5--6 and the near-resonant parents of families 7--10; see Section~\ref{subsec:lyapunov}). What differs is the parent stability type and proximity to resonance, not theoretical standing. `mode' is the centre frequency $\omega$; $N$ the number of corrected members; the $C_J$ and $T$ ranges are over the family; `max closure' is the largest periodicity residual $\|u(T)-u_0\|$ attained over the family's members (all $<10^{-10}$). Floquet class is by the stability index $\nu$ ($|\nu|<2$ stable).}
\label{tab:periodic_families}
\footnotesize
\setlength{\tabcolsep}{3.5pt}
\begin{tabular}{clcccccl}
\toprule
Fam & parent (arc, eq.) & $\omega$ & $N$ & $C_J$ range & $T$ range & max closure & Floquet / note \\
\midrule
1 & Case-1 stable, $(2.78,0.25)$ & $0.925$ & 22 & $[23.89,24.03]$ & $[6.78,6.79]$ & $5.3\times10^{-11}$ & Floquet stable \\
2 & Case-1 stable, $(2.78,0.25)$ & $0.406$ & 20 & $[24.03,24.04]$ & $[15.48,15.52]$ & $8.1\times10^{-11}$ & Floquet stable \\
3 & Case-2 stable, $(1.99,1.79)$ & $0.922$ & 23 & $[21.34,21.79]$ & $[6.68,6.82]$ & $3.5\times10^{-11}$ & Floquet stable \\
4 & Case-2 stable, $(1.99,1.79)$ & $0.390$ & 20 & $[21.79,21.80]$ & $[16.13,16.51]$ & $4.1\times10^{-11}$ & Floquet stable \\
5 & Case-1 sdl$\times$ctr, $(-0.97,0.51)$ & $3.738$ & 18 & $[13.33,13.43]$ & $[1.68,1.72]$ & $7.6\times10^{-11}$ & real unstable \\
6 & Case-2 sdl$\times$ctr, $(0.35,-0.64)$ & $7.435$ & 15 & $[24.28,24.29]$ & $[0.85,0.85]$ & $8.6\times10^{-11}$ & real unstable \\
7 & Case-1 $3{:}2$ reg., $(2.06,-0.76)$ & $0.851$ & 20 & $[15.23,15.25]$ & $[7.38,7.38]$ & $6.0\times10^{-11}$ & Floquet stable \\
8 & Case-1 $3{:}2$ reg., $(2.06,-0.76)$ & $0.564$ & 21 & $[15.25,15.29]$ & $[11.13,11.39]$ & $1.7\times10^{-11}$ & Floquet stable; $\nu\!\to\!-2$ \\
9 & Case-2 $5{:}2$ reg., $(-1.70,1.96)$ & $0.917$ & 21 & $[20.28,20.36]$ & $[6.79,6.85]$ & $9.0\times10^{-11}$ & Floquet stable \\
10 & Case-2 $5{:}2$ reg., $(-1.70,1.96)$ & $0.372$ & 18 & $[20.36,20.36]$ & $[16.89,17.17]$ & $2.6\times10^{-11}$ & Floquet stable; $\nu\!\to\!-2$ \\
\bottomrule
\end{tabular}
\end{table}



\subsection{Shooting formulation and linear seeding}
\label{subsec:periodic_shooting}

Write the rotating-frame state as
\[
\mathbf u=(x,y,\dot{x},\dot{y})^T
\]
and denote by $\phi_T(\mathbf u_0)$ the time-$T$ flow of
\eqref{eq:equations_of_motion}. A periodic orbit satisfies
\begin{equation}
\phi_T(\mathbf u_0)-\mathbf u_0=\mathbf 0.
\label{eq:periodicity_residual}
\end{equation}

Let
\[
L_j=(x^\ast,y^\ast,0,0)^T
\]
be a parent equilibrium with a simple centre eigenvalue
$+i\omega_k$ and corresponding complex eigenvector $\mathbf w_k$ of the
variational matrix \eqref{eq:variational_matrix}. The small-amplitude
linear motion provides the initial approximation
\[
\mathbf u(t)
=
L_j+
\varepsilon
\operatorname{Re}
\left(
e^{i\omega_k t}\mathbf w_k
\right).
\]
Using the real part of the selected eigenvector as the initial direction,
we take
\begin{equation}
\mathbf u_0^{(0)}
=
L_j+\varepsilon\widehat{\mathbf e}_k,
\qquad
\widehat{\mathbf e}_k
=
\frac{\operatorname{Re}\mathbf w_k}
{\|\operatorname{Re}\mathbf w_k\|},
\qquad
T_0=\frac{2\pi}{\omega_k}.
\label{eq:linear_periodic_seed}
\end{equation}
At a doubly-elliptic equilibrium the two simple centre modes are seeded
separately, producing two local Lyapunov families.

The periodicity equations alone do not isolate a unique member of a
one-parameter family because the autonomous flow also admits arbitrary
time translation along a periodic orbit. We therefore supplement the four
periodicity equations with a phase condition and an amplitude condition,
\begin{equation}
\widehat{\mathbf e}_k^T f(\mathbf u_0)=0,
\qquad
\widehat{\mathbf e}_k^T(\mathbf u_0-L_j)=a,
\label{eq:phase_amplitude_conditions}
\end{equation}
where $f$ denotes the rotating-frame vector field and $a$ is the prescribed
continuation amplitude.

The complete correction residual is consequently
\begin{equation}
\mathcal R(\mathbf u_0,T)=
\begin{pmatrix}
\phi_T(\mathbf u_0)-\mathbf u_0\\
\widehat{\mathbf e}_k^T f(\mathbf u_0)\\
\widehat{\mathbf e}_k^T(\mathbf u_0-L_j)-a
\end{pmatrix}.
\label{eq:shooting_residual}
\end{equation}
The state-transition matrix $\Phi(t)$ is integrated simultaneously from
\[
\dot{\Phi}=A(\mathbf u(t))\Phi,
\qquad
\Phi(0)=I,
\]
and the monodromy matrix is
\[
M=\Phi(T).
\]
The derivatives of the periodicity residual are
\[
D_{\mathbf u_0}
\bigl[\phi_T(\mathbf u_0)-\mathbf u_0\bigr]
=
M-I,
\qquad
D_T
\bigl[\phi_T(\mathbf u_0)-\mathbf u_0\bigr]
=
f(\phi_T(\mathbf u_0)).
\]
Thus the $6\times5$ correction Jacobian is
\begin{equation}
J=
\begin{pmatrix}
M-I & f(\phi_T(\mathbf u_0))\\
\widehat{\mathbf e}_k^T A(\mathbf u_0) & 0\\
\widehat{\mathbf e}_k^T & 0
\end{pmatrix},
\label{eq:shooting_jacobian}
\end{equation}
and each Newton correction is obtained from the least-squares problem
\[
\delta X
=
\arg\min_{\delta}
\|J\delta+\mathcal R\|_2.
\]
The phase and amplitude constraints remove the two neutral directions
associated with time translation and displacement along the local
one-parameter family. The resulting overdetermined formulation retains all
four periodicity equations symmetrically rather than eliminating a state
coordinate through the Jacobi integral.


\subsection{Continuation and numerical validation}
\label{subsec:periodic_continuation}

The state, period, and state-transition matrix are integrated with DOP853
using
\[
\mathrm{rtol}=\mathrm{atol}=10^{-12}.
\]
Newton corrections are capped according to
\[
\|\delta\mathbf u_0\|\leq0.5,
\qquad
|\delta T|\leq0.5T,
\]
with $T>0$ enforced throughout.

A corrected orbit is accepted only when
\begin{equation}
\|\phi_T(\mathbf u_0)-\mathbf u_0\|_2<10^{-10},
\qquad
|\widehat{\mathbf e}_k^T f(\mathbf u_0)|<10^{-10},
\qquad
|\widehat{\mathbf e}_k^T(\mathbf u_0-L_j)-a|<10^{-12}.
\label{eq:periodic_acceptance}
\end{equation}
Across the 198 accepted periodic solutions the observed maxima are
\[
\|\phi_T(\mathbf u_0)-\mathbf u_0\|_2
\leq9.0\times10^{-11},
\]
\[
|\widehat{\mathbf e}_k^T f(\mathbf u_0)|
\leq4.7\times10^{-12},
\qquad
|\widehat{\mathbf e}_k^T(\mathbf u_0-L_j)-a|
\leq3.9\times10^{-16}.
\]
The periodicity residual therefore provides a direct numerical certificate
that every orbit reported below closes to the stated tolerance.

We continue each family using the natural amplitude parameter
\[
a=
\widehat{\mathbf e}_k^T(\mathbf u_0-L_j),
\]
with the prescribed amplitudes forming the geometric sequence
\begin{equation}
a_0=2\times10^{-3},
\qquad
a_{n+1}=1.35\,a_n.
\label{eq:amplitude_sequence}
\end{equation}
Each new member is predicted from its predecessor and then corrected by the
shooting system above, with at most 40 Newton iterations. For geometric
interpretation we also record
\[
A=
\max_t
\sqrt{(x(t)-x^\ast)^2+(y(t)-y^\ast)^2}.
\]

Several safeguards prevent continuation onto spurious or unrelated
solutions. Correction is terminated after an integration failure caused by
a near collision; solutions with
\[
T<0.4T_0
\]
are rejected to prevent collapse toward the trivial $T\rightarrow0$
solution; and the family is terminated if the geometric amplitude more than
doubles between consecutive members. Continuation also stops when the orbit
approaches a primary within $0.05$ or when
\[
A>2.5.
\]

Because the continuation parameter is the prescribed amplitude $a$, these
computations determine only the portion of each family reachable while $a$
remains a useful monotone local coordinate. We therefore make no claim that
the computed curves are maximal periodic-orbit families. Continuation
through a turning point in $a$ would require pseudo-arclength continuation
in the extended periodic-orbit system, which is not attempted here.

The ten parent modes were chosen to sample distinct dynamical situations.
Families~1--4 use both centre modes of one doubly-elliptic equilibrium on
each admissible arc. Families~5--6 use centre modes of selected
saddle$\times$centre equilibria. Families~7--8 are seeded from the
Case-1 equilibrium near the $3{:}2$ near-commensurability identified in
Section~\ref{subsec:centrefreq}, and families~9--10 from the distinct
Case-2 equilibrium near the $5{:}2$ near-commensurability. As established
in Section~\ref{subsec:centrefreq}, the latter parent is
\[
(x,y)\simeq(-1.70,1.96),
\qquad
\omega_1/\omega_2\simeq2.465,
\]
and is not the separate Case-2 equilibrium lying near the $2{:}1$ ratio.

For every family the smallest corrected members approach the intended
linear mode:
\[
\mathbf u_0(a)\rightarrow L_j,
\qquad
T(a)\rightarrow\frac{2\pi}{\omega_k}
\qquad\text{as }a\rightarrow0.
\]
For example,
\[
\frac{2\pi}{0.925}=6.793,\qquad
\frac{2\pi}{0.406}=15.48,\qquad
\frac{2\pi}{0.922}=6.815,
\]
in agreement with the smallest computed periods of families~1--3 to the
reported precision.


\subsection{Floquet stability}
\label{subsec:floquet}

The linear stability of each corrected periodic orbit is determined from
its monodromy matrix $M=\Phi(T)$. The Hamiltonian structure of the
rotating-frame problem implies reciprocal symmetry of the Floquet
multipliers. For a periodic orbit of this two-degree-of-freedom autonomous
Hamiltonian system there is also a trivial multiplier pair at $+1$,
associated with the flow and energy-family directions. Numerically this
pair is recovered to within a few times $10^{-6}$.

After removal of the trivial pair, one nontrivial reciprocal multiplier
pair remains. If its members are denoted $\lambda$ and $\lambda^{-1}$, we
use the stability index
\begin{equation}
\nu
=
\lambda+\lambda^{-1}.
\label{eq:floquet_index}
\end{equation}
For the planar problem $\nu$ is real. The nontrivial multipliers lie on the
unit circle when
\[
|\nu|<2
\]
and form a real reciprocal pair when
\[
|\nu|>2.
\]
The limiting values
\[
\nu=+2
\qquad\text{and}\qquad
\nu=-2
\]
correspond respectively to a $+1$ collision and a period-doubling
($-1$) collision of the nontrivial multipliers.

We use ``Floquet stable'' only to mean that the nontrivial multiplier pair
lies on the unit circle to the stated numerical tolerance; it does not
imply nonlinear stability of the periodic orbit. Since only one
nontrivial reciprocal pair remains in the planar problem, a Hamiltonian
collision between two distinct nontrivial multiplier pairs cannot occur;
the relevant local stability boundaries are therefore the values
$\nu=\pm2$.

Although DOP853 is not a symplectic integrator, the computed monodromy
matrices preserve the expected Hamiltonian structure to good numerical
accuracy. Over the 198 accepted orbits,
\[
|\det M-1|\leq1.5\times10^{-8},
\qquad
|\lambda_1\lambda_2-1|\leq1.5\times10^{-8},
\]
with median values near $10^{-13}$, while the trivial pair remains within
$5\times10^{-6}$ of $+1$ (median $\sim3\times10^{-8}$). These diagnostics
provide an independent numerical check on the Floquet classification.


\subsection{Computed families and near-period-doubling behaviour}
\label{subsec:periodic_results}

The ten families comprise 198 accepted periodic solutions. Their parent
equilibria, centre frequencies, Jacobi and period ranges, closure residuals,
and Floquet classifications are summarized in
Table~\ref{tab:periodic_families}.

Before turning to the individual families, it is useful to fix their
geometric character. Each parent equilibrium is a libration point of the
rotating four-primary configuration, and the associated Lyapunov orbits are
small oscillations \emph{about that point} rather than revolutions about any
individual primary. This is illustrated in Figure~\ref{fig:orbit_potential}.
For the Case-2 parent of family~3 the Hessian of $\Omega$ is positive
definite, with eigenvalues $0.0436$ and $2.9554$, so the parent is a strict
local minimum of the rotating-frame effective potential. It is also the
numerically identified global minimum of $2\Omega$ for this primary
configuration, with $2\Omega_{\min}\simeq21.7867$. The computed family
members satisfy $C_J\le2\Omega_{\min}$; consequently the corresponding
numerical Hill set contains no forbidden region and is the whole plane. The
nested family is therefore organized by the local potential minimum rather
than by confinement within a closed zero-velocity boundary.

\begin{figure}[htbp]
  \centering
  \includegraphics[width=0.8\textwidth]{periodic_orbit_potential.pdf}
  \caption{Effective-potential landscape for the Case-2 configuration
  supporting periodic-orbit family~3. Filled contours show $2\Omega(x,y)$,
  with the colour range clipped at large values for visualization;
  $2\Omega\to+\infty$ at the primary singularities and at large distance.
  Thin grey curves are representative level curves of $2\Omega$ at $22.29$,
  $23.29$, $25.29$, and $28.79$. The four primaries are marked by
  $m_0$--$m_3$, and the star denotes the doubly-elliptic parent equilibrium.
  The Hessian of $\Omega$ at this point has eigenvalues $0.0436$ and
  $2.9554$, so the parent is a strict local minimum of the effective
  potential. Successive members of the associated Lyapunov family are shown
  from small to larger amplitude, illustrating libration about the
  equilibrium of the rotating four-primary configuration rather than orbital
  motion about any individual primary.}
  \label{fig:orbit_potential}
\end{figure}

Families~1--4 originate from the selected doubly-elliptic equilibria on the
two admissible arcs. Both centre modes of each parent produce corrected
Lyapunov families whose nontrivial Floquet multipliers remain on the unit
circle throughout the computed intervals. Figures~\ref{fig:family_case1}
and~\ref{fig:family_case2} show representative small, intermediate, and
large members. The smallest solutions reproduce the corresponding linear
centre modes, while larger members deform substantially without losing
Floquet stability over the continuation range. Figure~\ref{fig:famevol}
shows the period, geometric amplitude, and Floquet index of all ten
families against $C_J$.

Families~5--6 demonstrate the contrasting behaviour of
saddle$\times$centre parents. Their selected centre modes satisfy the
Lyapunov-centre non-resonance condition discussed in
Section~\ref{subsec:centrefreq}, so genuine local periodic families exist
despite the hyperbolic companion direction. Their nontrivial Floquet
multipliers, however, form a strongly real reciprocal pair, giving
\[
|\nu|\sim10^3
\]
over the computed portions. These families are therefore real unstable.
A representative member is shown in Figure~\ref{fig:selected_orbits}.

Families~7--10 sample the near-commensurate doubly-elliptic regions.
Families~7--8 originate near the Case-1 $3{:}2$ ratio, while
families~9--10 originate near the Case-2 $5{:}2$ ratio. These ratios are
near-commensurabilities rather than exact resonances at the parent
equilibria, so their local existence as Lyapunov families is unaffected.

Two of these families approach the period-doubling boundary
$\nu=-2$ as their amplitudes increase. Fine continuation near the minimum
of
\[
|\nu+2|
\]
shows that family~8 remains approximately
\[
2\times10^{-3}
\]
away from the threshold, whereas family~10 approaches much more closely,
reaching
\begin{equation}
\min_a|\nu+2|
\simeq1.0\times10^{-7},
\qquad
\nu\simeq-1.99999990,
\qquad
A\simeq0.28.
\label{eq:near_period_doubling}
\end{equation}
The index subsequently moves away from $-2$ rather than crossing it
(Figure~\ref{fig:period_doubling_scan}).

The approach was re-examined with an ultra-fine amplitude sweep of $51$
samples spanning amplitudes $0.258$ to $0.298$ (step $8\times10^{-4}$) around
the minimum, which returns the same closest approach,
$\min_a|\nu+2|\simeq1.0\times10^{-7}$ at $\nu\simeq-1.99999990$, with the index
turning back rather than crossing. All orbits are integrated at
$\mathrm{rtol}=\mathrm{atol}=10^{-12}$
(Section~\ref{subsec:periodic_continuation}). Recomputing the entire sweep at
$\mathrm{rtol}=\mathrm{atol}=10^{-14}$, the double-precision limit, reproduces
the same minimum to within $2\times10^{-16}$ and the same turn-back, so the
near-degeneracy is not an integration artefact; the monodromy diagnostics of
Section~\ref{subsec:floquet} likewise remain reliable at this point. We therefore find no
period-doubling crossing over the computed continuation interval. In
particular, the calculation provides no evidence for a local period-doubled
branch bifurcating from family~10 within that interval.

A corresponding scan of the computed families for crossings of the
selected low-order multiplier conditions
\[
\nu=2,\qquad
\nu=0,\qquad
\nu=-1,\qquad
\nu=-2
\]
finds none over the continuation ranges investigated. This is a negative
result restricted to the computed families and amplitudes; it is not a
claim that secondary resonant or period-doubled periodic solutions cannot
exist elsewhere in phase or parameter space.

Finally, four distinct notions must remain separate. Spectral stability of
an equilibrium refers to the eigenvalues of
\eqref{eq:variational_matrix}; Floquet stability of a periodic orbit refers
to its monodromy multipliers; numerical boundedness over a finite
integration interval establishes neither; and nonlinear or KAM stability
is not proved in this work. Every periodic orbit reported here satisfies
the closure conditions \eqref{eq:periodic_acceptance}, and every periodic
stability label is a Floquet statement.

\begin{figure}[htbp]
  \centering
  \includegraphics[width=\textwidth]{periodic_family_case1_mode2.pdf}
  \caption{Stable Lyapunov family from the Case-1 arc equilibrium (mode
  $\omega_2$): small (linear-mode seed), moderate, and large members. Green
  $+$: parent equilibrium; black dot: initial point; arrow: direction; stars:
  primaries in view. Floquet stable throughout.}
  \label{fig:family_case1}
\end{figure}
\begin{figure}[htbp]
  \centering
  \includegraphics[width=\textwidth]{periodic_family_case2_mode1.pdf}
  \caption{Stable Lyapunov family from the Case-2 arc equilibrium (mode
  $\omega_1$). Floquet stable throughout.}
  \label{fig:family_case2}
\end{figure}
\begin{figure}[htbp]
  \centering
  \includegraphics[width=\textwidth]{family_evolution.pdf}
  \caption{Family evolution: period $T$, geometric amplitude $A$, and Floquet
  index $\nu$ (symlog) against $C_J$. Real-unstable families (5, 6) sit at
  $\nu\sim10^{3}$; all others have $|\nu|<2$ over their computed range.}
  \label{fig:famevol}
\end{figure}
\begin{figure}[htbp]
  \centering
  \includegraphics[width=\textwidth]{selected_resonant_orbits.pdf}
  \caption{Representative periodic orbits: a real-unstable orbit from a
  saddle$\times$centre parent (family~5), a near-period-doubling member of
  the Case-2 $5{:}2$ family (family~10), a multi-lobed member of the
  Case-1 $3{:}2$ family (family~8), and a large Floquet-stable orbit
  (family~2). Titles give the geometric amplitude, Floquet index $\nu$,
  and Floquet class.}
  \label{fig:selected_orbits}
\end{figure}
\begin{figure}[htbp]
  \centering
  \includegraphics[width=0.72\textwidth]{period_doubling_scan.pdf}
  \caption{Fine scan of the Floquet index toward $\nu=-2$. Family~10 approaches
  the threshold to within $\sim10^{-7}$ and returns; the multiplier does not
  cross $-1$. No period-doubling bifurcation occurs.}
  \label{fig:period_doubling_scan}
\end{figure}
